\documentclass[11pt,a4paper]{article}
% Preambulo comun de los reportes de Tecnologias Celulares.
% Cada main.tex hace: \documentclass ... % Preambulo comun de los reportes de Tecnologias Celulares.
% Cada main.tex hace: \documentclass ... % Preambulo comun de los reportes de Tecnologias Celulares.
% Cada main.tex hace: \documentclass ... \input{../comun/preambulo.tex}
\usepackage[utf8]{inputenc}
\usepackage[T1]{fontenc}
\usepackage[spanish,mexico]{babel}
\usepackage{graphicx}
\usepackage[margin=2.5cm]{geometry}
\usepackage{amsmath}
\usepackage{amssymb}
\usepackage{booktabs}
\usepackage{tabularx}
\usepackage{array}
\usepackage{float}
\usepackage{setspace}
\usepackage{parskip}
\usepackage{fancyhdr}
\usepackage{textcomp}
\usepackage{xcolor}
\usepackage{caption}
\usepackage{fvextra}
\usepackage[hidelinks]{hyperref}

\onehalfspacing
\setlength{\parskip}{6pt}
\sloppy
\setlength{\emergencystretch}{3em}
\setlength{\headheight}{40pt}
\captionsetup{font=small,labelfont=bf}
\setlength{\tabcolsep}{6pt}

\pagestyle{fancy}
\fancyhf{}
\fancyhead[L]{\includegraphics[width=1cm]{/home/yovick/repos/UAZ-IEI/Assets/UIE.png}}
\fancyhead[R]{\small Tecnologías Celulares}
\fancyfoot[C]{\thepage}
\renewcommand{\headrulewidth}{0.4pt}

\DefineVerbatimEnvironment{codigo}{Verbatim}{%
  breaklines=true,breakanywhere=true,fontsize=\small,frame=single,
  framesep=3pt,rulecolor=\color{gray!50}}

% Figura centrada con pie: \figura[ancho]{archivo}{leyenda}
\newcommand{\figura}[3][0.85]{%
  \begin{center}
    \includegraphics[width=#1\linewidth]{#2}\\[2pt]
    {\small #3}
  \end{center}}

% Portada: \portada{numero}{titulo de la practica}
\newcommand{\portada}[2]{%
\begin{titlepage}
  \thispagestyle{empty}
  \begin{center}
    \vspace*{1cm}
    \includegraphics[width=4.5cm]{/home/yovick/repos/UAZ-IEI/Assets/uaz.png}
    \vspace{0.8cm}

    {\large Universidad Autónoma de Zacatecas}\\[4pt]
    {\large Unidad Académica de Ingeniería Eléctrica}\\[4pt]
    {\normalsize Ingeniería Electrónica Industrial con especialidad en Telecomunicaciones}

    \vspace{1.8cm}
    {\LARGE\bfseries Práctica #1}\\[10pt]
    {\Large\bfseries #2}
  \end{center}

  \vspace{2.2cm}
  \begin{center}
    \begin{tabular}{rl}
      \textbf{Materia:} & Tecnologías Celulares\\
      \textbf{Docente:} & Dr. Víktor Iván Rodríguez Abdalá\\
      \textbf{Alumno:}  & Arnold Yovick Ríos Zaldívar\\
      \textbf{Fecha:}   & Noviembre de 2026\\
    \end{tabular}
  \end{center}
  \vfill
\end{titlepage}}

\usepackage[utf8]{inputenc}
\usepackage[T1]{fontenc}
\usepackage[spanish,mexico]{babel}
\usepackage{graphicx}
\usepackage[margin=2.5cm]{geometry}
\usepackage{amsmath}
\usepackage{amssymb}
\usepackage{booktabs}
\usepackage{tabularx}
\usepackage{array}
\usepackage{float}
\usepackage{setspace}
\usepackage{parskip}
\usepackage{fancyhdr}
\usepackage{textcomp}
\usepackage{xcolor}
\usepackage{caption}
\usepackage{fvextra}
\usepackage[hidelinks]{hyperref}

\onehalfspacing
\setlength{\parskip}{6pt}
\sloppy
\setlength{\emergencystretch}{3em}
\setlength{\headheight}{40pt}
\captionsetup{font=small,labelfont=bf}
\setlength{\tabcolsep}{6pt}

\pagestyle{fancy}
\fancyhf{}
\fancyhead[L]{\includegraphics[width=1cm]{/home/yovick/repos/UAZ-IEI/Assets/UIE.png}}
\fancyhead[R]{\small Tecnologías Celulares}
\fancyfoot[C]{\thepage}
\renewcommand{\headrulewidth}{0.4pt}

\DefineVerbatimEnvironment{codigo}{Verbatim}{%
  breaklines=true,breakanywhere=true,fontsize=\small,frame=single,
  framesep=3pt,rulecolor=\color{gray!50}}

% Figura centrada con pie: \figura[ancho]{archivo}{leyenda}
\newcommand{\figura}[3][0.85]{%
  \begin{center}
    \includegraphics[width=#1\linewidth]{#2}\\[2pt]
    {\small #3}
  \end{center}}

% Portada: \portada{numero}{titulo de la practica}
\newcommand{\portada}[2]{%
\begin{titlepage}
  \thispagestyle{empty}
  \begin{center}
    \vspace*{1cm}
    \includegraphics[width=4.5cm]{/home/yovick/repos/UAZ-IEI/Assets/uaz.png}
    \vspace{0.8cm}

    {\large Universidad Autónoma de Zacatecas}\\[4pt]
    {\large Unidad Académica de Ingeniería Eléctrica}\\[4pt]
    {\normalsize Ingeniería Electrónica Industrial con especialidad en Telecomunicaciones}

    \vspace{1.8cm}
    {\LARGE\bfseries Práctica #1}\\[10pt]
    {\Large\bfseries #2}
  \end{center}

  \vspace{2.2cm}
  \begin{center}
    \begin{tabular}{rl}
      \textbf{Materia:} & Tecnologías Celulares\\
      \textbf{Docente:} & Dr. Víktor Iván Rodríguez Abdalá\\
      \textbf{Alumno:}  & Arnold Yovick Ríos Zaldívar\\
      \textbf{Fecha:}   & Noviembre de 2026\\
    \end{tabular}
  \end{center}
  \vfill
\end{titlepage}}

\usepackage[utf8]{inputenc}
\usepackage[T1]{fontenc}
\usepackage[spanish,mexico]{babel}
\usepackage{graphicx}
\usepackage[margin=2.5cm]{geometry}
\usepackage{amsmath}
\usepackage{amssymb}
\usepackage{booktabs}
\usepackage{tabularx}
\usepackage{array}
\usepackage{float}
\usepackage{setspace}
\usepackage{parskip}
\usepackage{fancyhdr}
\usepackage{textcomp}
\usepackage{xcolor}
\usepackage{caption}
\usepackage{fvextra}
\usepackage[hidelinks]{hyperref}

\onehalfspacing
\setlength{\parskip}{6pt}
\sloppy
\setlength{\emergencystretch}{3em}
\setlength{\headheight}{40pt}
\captionsetup{font=small,labelfont=bf}
\setlength{\tabcolsep}{6pt}

\pagestyle{fancy}
\fancyhf{}
\fancyhead[L]{\includegraphics[width=1cm]{/home/yovick/repos/UAZ-IEI/Assets/UIE.png}}
\fancyhead[R]{\small Tecnologías Celulares}
\fancyfoot[C]{\thepage}
\renewcommand{\headrulewidth}{0.4pt}

\DefineVerbatimEnvironment{codigo}{Verbatim}{%
  breaklines=true,breakanywhere=true,fontsize=\small,frame=single,
  framesep=3pt,rulecolor=\color{gray!50}}

% Figura centrada con pie: \figura[ancho]{archivo}{leyenda}
\newcommand{\figura}[3][0.85]{%
  \begin{center}
    \includegraphics[width=#1\linewidth]{#2}\\[2pt]
    {\small #3}
  \end{center}}

% Portada: \portada{numero}{titulo de la practica}
\newcommand{\portada}[2]{%
\begin{titlepage}
  \thispagestyle{empty}
  \begin{center}
    \vspace*{1cm}
    \includegraphics[width=4.5cm]{/home/yovick/repos/UAZ-IEI/Assets/uaz.png}
    \vspace{0.8cm}

    {\large Universidad Autónoma de Zacatecas}\\[4pt]
    {\large Unidad Académica de Ingeniería Eléctrica}\\[4pt]
    {\normalsize Ingeniería Electrónica Industrial con especialidad en Telecomunicaciones}

    \vspace{1.8cm}
    {\LARGE\bfseries Práctica #1}\\[10pt]
    {\Large\bfseries #2}
  \end{center}

  \vspace{2.2cm}
  \begin{center}
    \begin{tabular}{rl}
      \textbf{Materia:} & Tecnologías Celulares\\
      \textbf{Docente:} & Dr. Víktor Iván Rodríguez Abdalá\\
      \textbf{Alumno:}  & Arnold Yovick Ríos Zaldívar\\
      \textbf{Fecha:}   & Noviembre de 2026\\
    \end{tabular}
  \end{center}
  \vfill
\end{titlepage}}


\begin{document}

\portada{3}{Implementación de un sistema DSSS}

\textbf{Script:} \texttt{practica3.py} (Python con NumPy/SciPy)

\section{Objetivo}

Implementar un sistema de espectro disperso por secuencia directa (DSSS), verificar el
ensanchamiento espectral que produce el código PN y evaluar la recuperación de los datos
frente al ruido mediante la tasa de error de bit.

\section{Fundamento teórico}

En DSSS la señal de información $b(t)$ se multiplica por un código pseudoaleatorio $c(t)$
de velocidad mucho mayor:
\begin{equation}
  s(t) = b(t)\,c(t).
\end{equation}
Cada bit de datos se transmite como $N$ chips. Si $R_b$ es la velocidad de datos y $R_c$
la velocidad de chips, la ganancia de procesamiento es
\begin{equation}
  G_p = \frac{R_c}{R_b} = N,
  \qquad
  G_{p,\mathrm{dB}} = 10\log_{10} N.
\end{equation}
Para $N=31$, $G_p=14.91$ dB. El ancho de banda ocupado crece en el mismo factor, de modo
que la densidad espectral de potencia baja proporcionalmente.

En el receptor, la correlación con el código local,
\begin{equation}
  z = \sum_{k=0}^{N-1} r[k]\,c[k],
\end{equation}
suma coherentemente la señal útil (ganancia $N$) y solo suma incoherentemente el ruido
(ganancia $\sqrt{N}$), con lo que la relación señal a ruido de la decisión mejora en $N$.
Con chips de energía unitaria, la energía por bit es $E_b=N$ y la BER coincide con la de
BPSK:
\begin{equation}
  P_b = Q\!\left(\sqrt{\frac{2E_b}{N_0}}\right).
\end{equation}
La ventaja del DSSS no es reducir la BER para un $E_b/N_0$ dado, sino tolerar
interferencia de banda angosta: al desdispersar, una interferencia concentrada en
frecuencia se reparte y se atenúa en $G_p$.

\section{Desarrollo experimental}

\subsection{Parámetros}

\begin{tabularx}{\linewidth}{@{}*{2}{>{\raggedright\arraybackslash}X}@{}}
  \toprule
  Parámetro & Valor \\
  \midrule
  Bits de la demostración & 100 \\
  Código PN & $m=5$, $x^5+x^2+1$, 31 chips \\
  Chips por bit & 31 \\
  Ganancia de procesamiento & 14.91 dB \\
  Modulación & BPSK bipolar \\
  Canal & AWGN complejo \\
  \bottomrule
\end{tabularx}

\subsection{Procedimiento}

\begin{enumerate}
  \item Se generan 100 bits aleatorios y se convierten a $\pm1$.
  \item Cada bit se multiplica por el código PN de 31 chips (\texttt{kron}), produciendo
        3100 chips.
  \item Se compara la PSD de la señal de datos (cada bit repetido durante 31 chips) con la
        de la señal dispersa.
  \item Para la curva de BER se simulan 200\,000 bits por punto de $E_b/N_0$. El ruido
        complejo tiene varianza $N/E_bN_0$ para que la energía por bit sea $E_b=N$ y la
        comparación con BPSK sea directa.
  \item El receptor correlaciona cada bloque de 31 chips con el código y decide por el
        signo de la parte real.
\end{enumerate}

\section{Código}

\begin{codigo}
def dispersar(bits, codigo):
    return np.kron(2 * np.asarray(bits) - 1, codigo)

def simular_ber(ebn0_db, codigo, nbits=200_000, rng=None):
    ebn0 = 10 ** (ebn0_db / 10)
    bits = rng.integers(0, 2, nbits)
    tx = dispersar(bits, codigo)
    sigma2 = N / ebn0                       # Eb = N chips de energia unitaria
    ruido = np.sqrt(sigma2 / 2) * (rng.standard_normal(tx.size)
                                   + 1j * rng.standard_normal(tx.size))
    z = np.real(desdispersar(tx + ruido, codigo, nbits))
    return float(np.mean((z > 0).astype(int) != bits))
\end{codigo}

\section{Resultados}

\subsection{Señal dispersa}

\figura{figs/senal_dsss.png}{Dos bits dispersados con el código PN de 31 chips.}

Cada bit se convierte en una ráfaga de 31 chips de $\pm1$; los cambios de bit se ven como
un cambio de polaridad del código, no como un pulso rectangular.

\subsection{Espectro}

\figura{figs/espectro.png}{PSD de la señal de datos y de la señal DSSS, normalizadas.}

La señal de datos concentra su energía en un lóbulo de ancho relativo
$1/31\approx0.032$; la señal dispersa ocupa todo el ancho de banda disponible
($f/F_s$ de $-0.5$ a $0.5$). El ensanchamiento es de 31 veces, es decir, 14.91 dB, y la
densidad espectral baja en la misma proporción.

\subsection{Tasa de error de bit}

\begin{tabular}{@{}cccc@{}}
  \toprule
  $E_b/N_0$ (dB) & BER simulada & BER teórica BPSK & Error relativo \\
  \midrule
  0  & 0.078995 & 0.078650 & 0.4\,\% \\
  2  & 0.038355 & 0.037506 & 2.3\,\% \\
  4  & 0.012360 & 0.012501 & 1.1\,\% \\
  6  & 0.002400 & 0.002388 & 0.5\,\% \\
  8  & 0.000170 & 0.000191 & 11\,\% \\
  10 & $<5\times10^{-6}$ & $3.87\times10^{-6}$ & --- \\
  \bottomrule
\end{tabular}

\figura{figs/ber.png}{BER del DSSS comparada con la curva teórica de BPSK.}

La simulación sigue la curva teórica dentro del error estadístico de Monte Carlo. A
10 dB no se registró ningún error en 200\,000 bits, lo que fija una cota superior de
$5\times10^{-6}$, consistente con el valor teórico $3.9\times10^{-6}$. A 8 dB la
discrepancia es mayor porque se esperan solo 38 errores en 200\,000 bits.

\section{Análisis de resultados}

\begin{itemize}
  \item \textbf{¿Cómo cambia el espectro después del proceso de dispersión?} Se ensancha
        por un factor $N=31$ y su densidad espectral baja 14.91 dB. La forma es
        aproximadamente plana porque el código PN tiene espectro plano.
  \item \textbf{¿Cuál es la función del código PN?} Multiplicar los datos por una secuencia
        de alta velocidad conocida por el receptor. Eso reparte la energía en banda ancha
        y permite recuperar los datos por correlación, rechazando interferencias y
        separando usuarios.
  \item \textbf{¿Por qué se recuperan los datos mediante correlación?} Porque el código
        local solo coincide con la réplica recibida en el retardo correcto; en ese punto
        las 31 contribuciones se suman coherentemente y producen un valor $\pm N$, mientras
        que el ruido y las señales no sincronizadas se suman de forma incoherente.
  \item \textbf{¿Qué ocurre cuando aumenta la longitud del código?} Aumenta la ganancia de
        procesamiento, el ancho de banda ocupado y la inmunidad a interferencia, pero se
        reduce la velocidad de datos para un ancho de banda fijo y crece la complejidad de
        sincronización.
\end{itemize}

\section{Cuestionario}

\begin{enumerate}
  \item \textbf{¿Qué es DSSS?} Espectro disperso por secuencia directa: una técnica en la
        que la información se multiplica por un código pseudoaleatorio de alta velocidad,
        de modo que la señal ocupa mucho más ancho de banda que el necesario para los
        datos.
  \item \textbf{¿Qué es la ganancia de procesamiento?} La relación entre la velocidad de
        chips y la de datos, $G_p=R_c/R_b=N$. Expresada en dB, es la mejora en la relación
        señal a ruido que se obtiene al desdispersar.
  \item \textbf{¿Por qué DSSS ofrece inmunidad al ruido?} Porque el receptor integra
        coherentemente $N$ chips de la señal útil, mientras que el ruido se integra
        incoherentemente: la SNR de la decisión mejora en $N$. Además, una interferencia de
        banda angosta se dispersa al desdispersar y se atenúa en $G_p$.
  \item \textbf{¿Qué aplicaciones tuvo DSSS en sistemas celulares?} Fue la base de IS-95
        (cdmaOne) y de UMTS/WCDMA, donde los códigos PN y las secuencias de
        scrambling separan usuarios y celdas. También se usa en Wi-Fi (DSSS en 802.11b),
        GPS y comunicaciones militares.
  \item \textbf{¿Cuál es la diferencia entre DSSS y FHSS?} En DSSS la señal ocupa
        simultáneamente toda la banda y se dispersa con un código; en FHSS la portadora
        salta de frecuencia en frecuencia siguiendo una secuencia pseudoaleatoria. DSSS
        dispersa en frecuencia, FHSS diversifica en frecuencia.
\end{enumerate}

\section{Conclusiones}

Se implementó un sistema DSSS con un código PN de 31 chips y se verificó un ensanchamiento
espectral de 31 veces (14.91 dB de ganancia de procesamiento). La BER simulada coincidió
con la curva teórica de BPSK para $E_b/N_0$ de 0 a 8 dB, y a 10 dB no se registraron
errores en 200\,000 bits. El resultado confirma que el DSSS no mejora la BER en un canal
AWGN puro, sino que aporta inmunidad frente a interferencia y permite el acceso múltiple
por código.

\end{document}
